% GENERATED FILE. Edit canonical lessons, then run npm run generate:books.
% canonical-source-hash: fnv1a64:848c115a2c995e57
% canonical-lessons: HI-C479-muh, HI-C479-masura, HI-C479-cusna, HI-C479-cabana, HI-C479-thukna

\chapter{The Mouth}
\label{ch:hi-the-mouth}
\hlchaptermodality{479}

\begin{chapteropening}
\textbf{By the end of this chapter:} \emph{I can talk about chewing, swallowing and the mouth.}

Complete the last lesson of chapter 479: I can talk about chewing, swallowing and the mouth.
\end{chapteropening}

\section[\texorpdfstring{\hi{मुँह} (mũh)}{मुँह (mũh)}]{\hi{मुँह} (mũh) — the mouth; the face}

\label{lesson:HI-C479-muh}



\begin{quote}
Before the new one: say the Hindi for a rib, then the Hindi for a kidney.
\end{quote}

\subsection*{You'll want to know: \hi{मुँह}}
\textbf{\hi{मुँह}} — \emph{mũh} — \textquotedblleft{}the mouth; the face\textquotedblright{}.

\begin{quote}
\hi{खाने} \hi{से} \hi{पहले} \hi{मुँह} \hi{और} \hi{हाथ} \hi{धो} \hi{लो।}
\end{quote}

\begin{quote}
\emph{khāne se pahle mũh aur hāth dho lo.}
\end{quote}

\begin{quote}
Wash your face and hands before eating.
\end{quote}

Masculine.

\begin{cousinweb}
\textbf{\hi{मुँह}} (\emph{mũh}) is a tadbhava: CDIAL lists Hindi \emph{mũh}, \textquotedblleft{}mouth, face\textquotedblright{}, under Sanskrit \emph{múkha}, \textquotedblleft{}mouth, face\textquotedblright{} (CDIAL 10158); the \emph{kh} wore down to \emph{h} and left a nasal vowel.
\end{cousinweb}

\begin{grammarlens}[title={What you've built}]
One more word for describing what you lived and felt.
\end{grammarlens}

\subsection*{Your turn}
\begin{itemize}
\raggedright
  \item \emph{Say it:} \emph{mũh}
  \item \emph{Say it:} \emph{mũh}, once more
  \item \emph{Say it:} \emph{khāne se pahle mũh aur hāth dho lo.} — Wash your face and hands before eating.
  \item \emph{Recall:} say the Hindi for a rib, then the Hindi for a kidney, then say \emph{mũh} again
\end{itemize}

\begin{tcolorbox}[breakable,colback=teal!4,colframe=teal!35!black,title={Before you move on}]
What does \hi{मुँह} mean? (\textquotedblleft{}The mouth; the face\textquotedblright{}.) Say it once more.
\end{tcolorbox}
\section[\texorpdfstring{\hi{मसूड़ा} (masūṛā)}{मसूड़ा (masūṛā)}]{\hi{मसूड़ा} (masūṛā) — a gum (of the teeth)}

\label{lesson:HI-C479-masura}



\begin{quote}
Before the new one: say the Hindi for a kidney, then the Hindi for the mouth.
\end{quote}

\subsection*{You'll want to know: \hi{मसूड़ा}}
\textbf{\hi{मसूड़ा}} — \emph{masūṛā} — \textquotedblleft{}a gum (of the teeth)\textquotedblright{}.

\begin{quote}
\hi{मीरा} \hi{के} \hi{मसूड़ों} \hi{में} \hi{सूजन} \hi{है।}
\end{quote}

\begin{quote}
\emph{Mīrā ke masūṛoṁ meṁ sūjan hai.}
\end{quote}

\begin{quote}
Mira's gums are swollen.
\end{quote}

Masculine.

\begin{cousinweb}
\textbf{\hi{मसूड़ा}} (\emph{masūṛā}) is a tadbhava: CDIAL lists Hindi \emph{masūṛā}, \textquotedblleft{}the gums\textquotedblright{}, under \emph{*māṁsapuṭaka}, a form it reconstructs from Sanskrit \emph{māṁsa}, \textquotedblleft{}flesh\textquotedblright{} (CDIAL 9984).
\end{cousinweb}

\begin{grammarlens}[title={What you've built}]
One more word for describing what you lived and felt.
\end{grammarlens}

\subsection*{Your turn}
\begin{itemize}
\raggedright
  \item \emph{Say it:} \emph{masūṛā}
  \item \emph{Say it:} \emph{masūṛā}, once more
  \item \emph{Say it:} \emph{Mīrā ke masūṛoṁ meṁ sūjan hai.} — Mira's gums are swollen.
  \item \emph{Recall:} say the Hindi for a kidney, then the Hindi for the mouth, then say \emph{masūṛā} again
\end{itemize}

\begin{tcolorbox}[breakable,colback=teal!4,colframe=teal!35!black,title={Before you move on}]
What does \hi{मसूड़ा} mean? (\textquotedblleft{}A gum (of the teeth)\textquotedblright{}.) Say it once more.
\end{tcolorbox}
\section[\texorpdfstring{\hi{चूसना} (cūsnā)}{चूसना (cūsnā)}]{\hi{चूसना} (cūsnā) — to suck}

\label{lesson:HI-C479-cusna}



\begin{quote}
Before the new one: say the Hindi for the mouth, then the Hindi for a gum (of the teeth).
\end{quote}

\subsection*{You'll want to know: \hi{चूसना}}
\textbf{\hi{चूसना}} — \emph{cūsnā} — \textquotedblleft{}to suck\textquotedblright{}.

\begin{quote}
\hi{बच्चा} \hi{अपना} \hi{अंगूठा} \hi{चूसता} \hi{है।}
\end{quote}

\begin{quote}
\emph{baccā apnā aṅgūṭhā cūstā hai.}
\end{quote}

\begin{quote}
The baby sucks its thumb.
\end{quote}

Past, with \hi{ने} (he or she …): \hi{उसने} \hi{आम} \hi{चूसा}, \emph{usne ām cūsā}.

\begin{cousinweb}
\textbf{\hi{चूसना}} (\emph{cūsnā}) is a tadbhava: CDIAL lists Hindi \emph{cūsnā}, \textquotedblleft{}to suck\textquotedblright{}, under Sanskrit \emph{cūṣati}, \textquotedblleft{}sucks\textquotedblright{} (CDIAL 4898). A ripe mango is often sucked rather than cut.
\end{cousinweb}

\begin{grammarlens}[title={What you've built}]
One more word for describing what you lived and felt.
\end{grammarlens}

\subsection*{Your turn}
\begin{itemize}
\raggedright
  \item \emph{Say it:} \emph{cūsnā}
  \item \emph{Say it:} \emph{cūsnā}, once more
  \item \emph{Say it:} \emph{baccā apnā aṅgūṭhā cūstā hai.} — The baby sucks its thumb.
  \item \emph{Recall:} say the Hindi for the mouth, then the Hindi for a gum (of the teeth), then say \emph{cūsnā} again
\end{itemize}

\begin{tcolorbox}[breakable,colback=teal!4,colframe=teal!35!black,title={Before you move on}]
What does \hi{चूसना} mean? (\textquotedblleft{}To suck\textquotedblright{}.) Say it once more.
\end{tcolorbox}
\section[\texorpdfstring{\hi{चबाना} (cabānā)}{चबाना (cabānā)}]{\hi{चबाना} (cabānā) — to chew}

\label{lesson:HI-C479-cabana}



\begin{quote}
Before the new one: say the Hindi for a gum (of the teeth), then the Hindi for to suck.
\end{quote}

\subsection*{You'll want to know: \hi{चबाना}}
\textbf{\hi{चबाना}} — \emph{cabānā} — \textquotedblleft{}to chew\textquotedblright{}.

\begin{quote}
\hi{दादा} \hi{के} \hi{दाँत} \hi{कमज़ोर} \hi{हैं}, \hi{इसलिए} \hi{वह} \hi{रोटी} \hi{नहीं} \hi{चबा} \hi{सकते।}
\end{quote}

\begin{quote}
\emph{dādā ke dā̃t kamzor haiṁ, isliye vah roṭī nahīṃ cabā sakte.}
\end{quote}

\begin{quote}
Grandfather's teeth are weak, so he can't chew bread.
\end{quote}

Past, with \hi{ने} (he or she …): \hi{उसने} \hi{रोटी} \hi{चबाई}, \emph{usne roṭī cabāī}.

\begin{cousinweb}
\textbf{\hi{चबाना}} (\emph{cabānā}) is a tadbhava: CDIAL lists Hindi \emph{cabānā}, \textquotedblleft{}to bite the lip, gnaw, crunch\textquotedblright{}, under Sanskrit \emph{cárvati}, \textquotedblleft{}chews\textquotedblright{} (CDIAL 4711).
\end{cousinweb}

\begin{grammarlens}[title={What you've built}]
One more word for describing what you lived and felt.
\end{grammarlens}

\subsection*{Your turn}
\begin{itemize}
\raggedright
  \item \emph{Say it:} \emph{cabānā}
  \item \emph{Say it:} \emph{cabānā}, once more
  \item \emph{Say it:} \emph{dādā ke dā̃t kamzor haiṁ, isliye vah roṭī nahīṃ cabā sakte.} — Grandfather's teeth are weak, so he can't chew bread.
  \item \emph{Recall:} say the Hindi for a gum (of the teeth), then the Hindi for to suck, then say \emph{cabānā} again
\end{itemize}

\begin{tcolorbox}[breakable,colback=teal!4,colframe=teal!35!black,title={Before you move on}]
What does \hi{चबाना} mean? (\textquotedblleft{}To chew\textquotedblright{}.) Say it once more.
\end{tcolorbox}
\section[\texorpdfstring{\hi{थूकना} (thūknā)}{थूकना (thūknā)}]{\hi{थूकना} (thūknā) — to spit}

\label{lesson:HI-C479-thukna}



\begin{quote}
Before the new one: say the Hindi for to suck, then the Hindi for to chew.
\end{quote}

\subsection*{You'll want to know: \hi{थूकना}}
\textbf{\hi{थूकना}} — \emph{thūknā} — \textquotedblleft{}to spit\textquotedblright{}.

\begin{quote}
\hi{सड़क} \hi{पर} \hi{थूकना} \hi{बुरी} \hi{आदत} \hi{है।}
\end{quote}

\begin{quote}
\emph{saṛak par thūknā burī ādat hai.}
\end{quote}

\begin{quote}
Spitting in the street is a bad habit.
\end{quote}

Past, with \hi{ने} (he or she …): \hi{उसने} \hi{सड़क} \hi{पर} \hi{थूका}, \emph{usne saṛak par thūkā}.

\begin{cousinweb}
\textbf{\hi{थूकना}} (\emph{thūknā}) is a tadbhava: CDIAL lists Hindi \emph{thūknā}, \textquotedblleft{}spits\textquotedblright{}, under \emph{*thukk}, a form it reconstructs (CDIAL 6097); the word copies the sound it names.
\end{cousinweb}

\begin{grammarlens}[title={What you've built}]
One more word for describing what you lived and felt.
\end{grammarlens}

\subsection*{Your turn}
\begin{itemize}
\raggedright
  \item \emph{Say it:} \emph{thūknā}
  \item \emph{Say it:} \emph{thūknā}, once more
  \item \emph{Say it:} \emph{saṛak par thūknā burī ādat hai.} — Spitting in the street is a bad habit.
  \item \emph{Recall:} say the Hindi for to suck, then the Hindi for to chew, then say \emph{thūknā} again
\end{itemize}

\begin{tcolorbox}[breakable,colback=teal!4,colframe=teal!35!black,title={Before you move on}]
What does \hi{थूकना} mean? (\textquotedblleft{}To spit\textquotedblright{}.) Say it once more.
\end{tcolorbox}
